% GENERATED FILE. Edit canonical lessons, then run npm run generate:books.
% canonical-source-hash: fnv1a64:fb6efe5764d4aad5
% canonical-lessons: ES-C371-pequeno, ES-C371-gordo, ES-C371-joven, ES-C371-lento, ES-C371-dulce

\chapter{Small, Fat, Young, Slow, Sweet}
\label{ch:pequeno-y-gordo}
\hlchaptermodality{371}

\begin{chapteropening}
\textbf{By the end of this chapter:} \emph{I can say pequeño, gordo, joven, lento and dulce, and put five more judgements on a thing without reaching for a number.}

Complete ES-C371-dulce: put five judgements on a thing without reaching for a number.
\end{chapteropening}

\section[pequeño]{pequeño — the word with no ancestors}

\label{lesson:ES-C371-pequeno}



\begin{quote}
Say \emph{the cave}, then \emph{the coast}. (\emph{La cueva}; \emph{la costa}.)
\end{quote}

\subsection*{What to know first}
\begin{itemize}
\raggedright
  \item grande — the word this one answers.
  \item el tamaño — what both of them are about.
\end{itemize}

\begin{sounds}
\begin{itemize}
\raggedright
  \item \emph{pe-QUE-ño}, three beats with the weight on the second. The \emph{qu} is a hard \emph{k}, and the \emph{ñ} is the \emph{ny} of English \emph{canyon}. $\to$ pronunciation reference
\end{itemize}
\end{sounds}

\begin{cousinweb}
\textbf{pequeño} — \textquotedblleft{}small.\textquotedblright{}

Nearly every word in this book has a Latin ancestor standing behind it. This one does not, and the Spanish dictionary of record says so in as many words: an \textbf{expressive} word, common to all the Romance languages, with no etymon given.

What that means is worth unpacking, because it is not a confession of ignorance. Some words are not inherited. They are \textbf{made}, over and over, out of the sounds people reach for when they are talking to something small — a child, an animal, a thing they are fond of. Across the Romance languages those sounds cluster around a tight little syllable of the \emph{pikk-} or \emph{pitt-} shape.

Italian \textbf{piccolo}. French \textbf{petit}. Spanish \textbf{pequeño}. They are not three descendants of one Latin word. They are three separate helpings from the same nursery cupboard.

English then borrowed two of them and can now say the idea three ways: \textbf{small} from its own stock, \textbf{petite} from French, and the \textbf{piccolo}, the little flute, from Italian.

Take the lesson as well as the word. An etymology that ends in \emph{nobody knows} is sometimes a gap in the record, and sometimes it is the answer: this word was never handed down, so there is nothing upstream to find.
\end{cousinweb}

\begin{grammarlens}[title={What you've built}]
You can now set the two sizes against each other: \emph{grande} and \emph{pequeño}.
\end{grammarlens}

\subsection*{Your turn}
Say these aloud:

\begin{itemize}
\raggedright
  \item \emph{pequeño}
  \item \emph{muy pequeño}
  \item \emph{una casa pequeña}
\end{itemize}

\begin{tcolorbox}[breakable,colback=teal!4,colframe=teal!35!black,title={Before you move on}]
Does \emph{pequeño} have a Latin ancestor? (No — it is an expressive word with no etymon.) And which two Romance small-words did English borrow? (\emph{Petite} and \emph{piccolo}.)
\end{tcolorbox}
\section[gordo]{gordo — the word that may have gone to Rome and come back}

\label{lesson:ES-C371-gordo}



\begin{quote}
Say \emph{the coast}, then \emph{small}. (\emph{La costa}; \emph{pequeño}.)
\end{quote}

\subsection*{What to know first}
\begin{itemize}
\raggedright
  \item delgado / flaco — the two words this one answers.
  \item la carne — where the idea usually lands.
\end{itemize}

\begin{sounds}
\begin{itemize}
\raggedright
  \item \emph{GOR-do}, two beats with the weight on the first. The \emph{g} is hard, and the \emph{d} after \emph{r} is soft, closer to the \emph{th} of English \emph{this}. $\to$ pronunciation reference
\end{itemize}
\end{sounds}

\begin{cousinweb}
\textbf{gordo} — \textquotedblleft{}fat.\textquotedblright{}

Latin \textbf{gurdus} did not mean fat. It meant \textbf{dull, blunt, thick-witted} — an insult about a mind rather than a body.

It comes with a footnote from a Roman who noticed it. Quintilian, writing about words that had come into Latin from elsewhere, mentions \emph{gurdi}, which ordinary people say for stupid, and adds: \textbf{I have heard it took its origin from Spain.}

Read the verb carefully. He heard it. He is passing on a report, not making a finding, and nobody since has been able to prove or disprove him. But if he is right, then \emph{gordo} is a word that left the Iberian peninsula, worked in Rome for several centuries, drifted from a blunt mind to a blunt shape, and came back home in Spanish with a new job. A word that emigrated and returned.

French kept a piece of it: \textbf{gourd}, meaning numb or stiff — still the idea of something that does not move quickly.

English kept none of it, which matters because English has a word that looks like it does. A \textbf{gourd}, the hard-shelled vegetable, is a different word entirely, come through French from Latin \emph{cucurbita}. Same spelling as the French adjective, no relation to it, and no relation to \emph{gordo}.
\end{cousinweb}

\begin{grammarlens}[title={What you've built}]
You can now describe a build both ways: \emph{gordo} and \emph{delgado}.
\end{grammarlens}

\subsection*{Your turn}
Say these aloud:

\begin{itemize}
\raggedright
  \item \emph{gordo}
  \item \emph{muy gordo}
  \item \emph{un gato gordo}
\end{itemize}

\begin{tcolorbox}[breakable,colback=teal!4,colframe=teal!35!black,title={Before you move on}]
What did Latin \emph{gurdus} mean? (Dull, thick-witted.) And did Quintilian state the Spanish origin, or report it? (He reported hearing it.)
\end{tcolorbox}
\section[joven]{joven — young twice in one language}

\label{lesson:ES-C371-joven}



\begin{quote}
Say \emph{small}, then \emph{fat}. (\emph{Pequeño}; \emph{gordo}.)
\end{quote}

\subsection*{What to know first}
\begin{itemize}
\raggedright
  \item viejo — the word this one answers.
  \item la edad — what both of them measure.
\end{itemize}

\begin{sounds}
\begin{itemize}
\raggedright
  \item \emph{JO-ven}, two beats with the weight on the first. The \emph{j} scrapes at the back of the mouth, and the \emph{v} sounds like a soft \emph{b}. $\to$ pronunciation reference
\end{itemize}
\end{sounds}

\begin{cousinweb}
\textbf{joven} — \textquotedblleft{}young.\textquotedblright{}

Latin \textbf{iuvenis}, young. English has this word twice, and the two arrivals look nothing alike:

\begin{itemize}
\raggedright
  \item \textbf{young} and \textbf{youth} were never borrowed. They came down the Germanic line from the same ancient root that fed \emph{iuvenis}, changing at every step, which is why they have kept no letters in common with the Latin;
  \item \textbf{juvenile} and \textbf{junior} were taken off the page much later. \emph{Junior} is a comparative worn short — \emph{iūnior}, \textquotedblleft{}younger.\textquotedblright{} Every Junior in a family name is being called the younger one in Latin.
\end{itemize}

\textbf{Rejuvenate} is not a Latin word. English built it in the nineteenth century out of Latin parts, and Rome would not have recognised it.

Then there is \textbf{June}, and it is a good example of an answer that has to stay hedged. \emph{Iūnius} is probably named for \textbf{Juno}, and Juno's own name may itself mean \textquotedblleft{}the young one\textquotedblright{} from this very root — so June reaches \emph{joven} at one remove, through two probablys.

The Romans were no surer. Ovid put three goddesses in a room arguing over who the month belonged to, let each make her case, and pointedly refused to declare a winner. When a poet who lived there will not settle it, a modern reader should not pretend to.
\end{cousinweb}

\begin{grammarlens}[title={What you've built}]
You can now place someone on the scale: \emph{joven} or \emph{viejo}.
\end{grammarlens}

\subsection*{Your turn}
Say these aloud:

\begin{itemize}
\raggedright
  \item \emph{joven}
  \item \emph{muy joven}
  \item \emph{una mujer joven}
\end{itemize}

\begin{tcolorbox}[breakable,colback=teal!4,colframe=teal!35!black,title={Before you move on}]
Which English word is the inherited cousin of \emph{joven}? (\emph{Young}.) And what does \emph{junior} mean taken literally? (Younger.)
\end{tcolorbox}
\section[lento]{lento — the word that arrived by post}

\label{lesson:ES-C371-lento}



\begin{quote}
Say \emph{fat}, then \emph{young}. (\emph{Gordo}; \emph{joven}.)
\end{quote}

\subsection*{What to know first}
\begin{itemize}
\raggedright
  \item más despacio — what you already say when you need someone to slow down.
  \item el tiempo — what slowness is measured against.
\end{itemize}

\begin{sounds}
\begin{itemize}
\raggedright
  \item \emph{LEN-to}, two beats with the weight on the first. The \emph{l} is light and forward, and the \emph{t} touches the teeth. $\to$ pronunciation reference
\end{itemize}
\end{sounds}

\begin{cousinweb}
\textbf{lento} — \textquotedblleft{}slow.\textquotedblright{}

Latin \textbf{lentus} started somewhere else. It meant \textbf{pliant, tough, flexible} — a branch that bends without breaking, pitch that stretches when you pull it. Slowness came out of that: a thing that yields slowly, that will not be hurried.

There is a fingerprint on this word telling you it was \textbf{borrowed off the page} rather than passed down by mouth, and it is the \emph{-nt-}. Latin consonant clusters wear down when a word travels through centuries of speech, and this one has not worn at all. Spanish took \emph{lento} from written Latin in the fifteenth century.

The form that \textbf{did} come down by mouth also survives, and it went somewhere unexpected: \textbf{liento}, meaning damp. It grew out of the clammy, sticky end of \emph{lentus} rather than the slow end. One Latin adjective, two Spanish words, one written and one spoken.

English has the word through French as \textbf{relent} — to bend back, to soften under pressure — and built \textbf{relentless} on top of it. Someone relentless is someone with no bend in them.

And a false friend to leave alone. A \textbf{lentil} is not a slow bean. Latin \emph{lēns, lentis} is a separate word whose own origin nobody has traced, resembling \emph{lentus} by accident.
\end{cousinweb}

\begin{grammarlens}[title={What you've built}]
You can now ask for a change of pace with an adjective, not only a phrase: \emph{muy lento}.
\end{grammarlens}

\subsection*{Your turn}
Say these aloud:

\begin{itemize}
\raggedright
  \item \emph{lento}
  \item \emph{muy lento}
  \item \emph{el autobús es lento}
\end{itemize}

\begin{tcolorbox}[breakable,colback=teal!4,colframe=teal!35!black,title={Before you move on}]
What did \emph{lentus} mean before it meant slow? (Pliant, flexible.) And what is the Spanish word that came down by mouth instead? (\emph{Liento}, damp.)
\end{tcolorbox}
\section[dulce]{dulce — a resemblance nobody can explain}

\label{lesson:ES-C371-dulce}



\begin{quote}
Say \emph{young}, then \emph{slow}. (\emph{Joven}; \emph{lento}.)
\end{quote}

\subsection*{What to know first}
\begin{itemize}
\raggedright
  \item la miel — the oldest sweet thing there is.
  \item el azúcar — the one that replaced it.
\end{itemize}

\begin{sounds}
\begin{itemize}
\raggedright
  \item \emph{DUL-ce}, two beats with the weight on the first. The \emph{c} before \emph{e} is soft — a \emph{th} in much of Spain, an \emph{s} across the Americas. $\to$ pronunciation reference
\end{itemize}
\end{sounds}

\begin{cousinweb}
\textbf{dulce} — \textquotedblleft{}sweet.\textquotedblright{}

Latin \textbf{dulcis}. But the word Spanish actually inherited by mouth was \textbf{duz} — the \emph{l} had worn away, the way letters do. It came back later, pulled into place by educated speakers who could read the Latin and wanted the spelling honoured. The same pressure keeps the \emph{l} standing in \emph{alto}.

English has the family in two quiet places: \textbf{dulcet}, a sweet sound, through French; and the \textbf{dulcimer}, the struck-string instrument, whose name most likely began as a phrase meaning sweet song.

Now the real reason this word is here.

Greek has \textbf{glykys}, sweet, which is where \textbf{glucose} and \textbf{glycerine} come from. \emph{Dulcis} and \emph{glykys} mean the same thing and look tantalisingly alike — and no rule of sound change allows a Latin \emph{d-} to match a Greek \emph{gl-}. Reconstruction can force a shape that would produce both, but then the missing vowel in the Latin has no explanation, and specialists say so.

The suggestion that has the most support is that neither language inherited the word: both \textbf{borrowed} it, from some third language around the Mediterranean that left no other trace.

Keep this one as a specimen. Most etymologies in this book end in an answer. This one ends in two languages holding nearly the same word and nobody able to say how either of them got it.
\end{cousinweb}

\begin{grammarlens}[title={What you've built}]
You can now describe a taste: \emph{muy dulce}, \emph{un café dulce}.
\end{grammarlens}

\subsection*{Your turn}
Say these aloud:

\begin{itemize}
\raggedright
  \item \emph{dulce}
  \item \emph{muy dulce}
  \item \emph{un café dulce}
\end{itemize}

\begin{tcolorbox}[breakable,colback=teal!4,colframe=teal!35!black,title={Before you move on}]
What was the older Spanish form of \emph{dulce}? (\emph{Duz}.) And why is the Greek \emph{glykys} resemblance a puzzle? (No sound rule turns a Latin \emph{d-} into a Greek \emph{gl-}.)
\end{tcolorbox}
